\section{7. Conclusion}

This study examined heterogeneous UAV coordination when a relay-node failure reconfigures legal communication paths and task-support relations. DRTP-SG-MAPPO retained the Single-Graph MAPPO backbone, PPO objective, reward, and execution-time information boundary, while using bounded adaptive reweighting over six predefined failure/topology groups under a fixed nominal-condition anchor. Its matched UTR-SG-MAPPO control used the same parameter count, perturbation support, training budget, and evaluation protocol, isolating adaptive rather than uniform training-group weighting as the intended contrast.

In the prospective formal five-seed cohort, DRTP had positive paired effects on canonical F0, \texttt{J_pert,mean\texttt{, and \texttt{J_pert,worst\texttt{, with all five seeds positive on these primary endpoints. It also had a lower mean timeout rate but a slightly higher collision rate, so the formal result is a mission-score improvement with a safety trade-off rather than evidence of uniform safety improvement. A post hoc evaluation on six training-excluded onset--duration members retained the positive direction in the formal cohort. MAPPO-NoGraph provided external task context but, because of its different inputs and capacity, did not replace the matched UTR--DRTP comparison.

The independent three-method cohort produced a different DRTP--UTR direction and included a catastrophic DRTP training seed. Cross-evaluation-tape and additional unseen-member diagnostics retained these cohort-specific directions, making a simple evaluation-tape explanation insufficient. DRTP is thus positioned as a bounded training-time strategy with demonstrated formal-cohort benefit under the frozen contract and an explicit cross-cohort reliability boundary. This combination of controlled task semantics, matched evidence, and transparent reliability reporting provides a rigorous basis for future mechanism-focused studies of adaptive MARL training reliability.